% ─────────────────────────────────────────────────────────────────────────────
% Sec. 1  Introduction. Numbers match Sec. 5 (fig_scaling, tab_main, tab_memcap), supp_scale (tab_scale) and the
% V* bbox statistics (80% of targets < 1 patch in the 448-px thumbnail; median 0.08% area, 1 of 24 tiles).
% ─────────────────────────────────────────────────────────────────────────────
\section{Introduction}
\label{sec:intro}

Memory is the main constraint in serving large language models. The KV cache grows with every token
in the context, and it, not compute, limits how many requests a GPU can hold at
once~\cite{pagedattention}, which has motivated a line of work on compressing it~\cite{h2o,streamingllm}.
Vision-language models (VLMs) make this problem much worse. To read small text and recognize small
objects, current VLMs split a high-resolution image into many tiles at near-native
resolution~\cite{internvl15,llavanext,internvl35}, and each tile becomes hundreds of visual tokens. Kept at
the native tile resolution of InternVL3.5-8B~\cite{internvl35}, an 8K image needs 188 tiles, \ie 48K
visual tokens. Its prefill
needs 12.9\,GB of GPU memory on top of the model weights and 17\,s to the first token, so a 48\,GB GPU
can hold only three such requests. Visual tokens, not text, now dominate the memory of a VLM, and their
number grows linearly with the image area.

Visual token pruning reduces the number of tokens that reach the language model, but it acts too late.
VisionZip~\cite{visionzip} and HiRED~\cite{hired} prune after the vision encoder, so every tile is still encoded.
FastV~\cite{fastv}, SparseVLM~\cite{sparsevlm}, and IF-Prune~\cite{ifprune} prune inside the language
model, so all tokens still pass through its first layers, and the peak memory of prefill is still set
by the full image. At 188 tiles these methods need more than 8\,GB, although they keep only 1{,}536
tokens. The obvious alternative, downscaling the image before encoding, loses exactly
the details that high resolution was meant to provide: on V$^*$~Bench~\cite{vstar}, 80\% of the target
objects become smaller than a single ViT patch in the global thumbnail. Local crops alone, on the other
hand, lose the layout and context of the whole image. Yet the answer usually lies in a small part of the
image: on V$^*$~Bench, the target objects cover a median of 0.08\% of the image and fall into a single one
of 24 tiles, while token pruning encodes all of them.

The question is therefore not how many visual tokens to keep, but which regions to compute at all. To
bound memory without losing these details, this decision must be made \emph{before} the regions are
encoded (\cref{fig:method}). This is feasible: low-resolution attention alone places the target within the top four of a median of 24 tiles on 67.0\% of the V$^*$~Bench questions, against 14.8\% for random selection
(\cref{sec:reliability}). We propose \method, a training-free method that follows the principle
\emph{look before you encode}. Like foveated vision~\cite{ram,glancefocus}, it combines a coarse view of
the whole scene with a few high-resolution glimpses: it spends half of the token budget on
low-resolution \emph{coverage tiles} of the whole image. A single partial forward pass of the VLM over
these tiles and the question yields a question-aware prior over the high-resolution \emph{detail tiles};
only the tiles at the selection boundary are re-scored with question attention. The selected detail
tiles are then encoded and streamed into the KV cache, and all other tiles never reach the vision
encoder. As a result, the vision encoder, the language model, and the KV cache all process a number of
tokens set by the budget, independently of the input resolution, and no model other than the VLM
itself is needed.

Our contributions are as follows.
\begin{itemize}
  \item We identify that existing visual token pruning cannot bound the memory of high-resolution VLMs,
        because it removes tokens only after they have been computed, and we propose to select image
        regions before encoding instead. We show on V$^*$~Bench that low-resolution attention can filter
        regions by question, and that re-scoring a few tiles at high resolution recovers most of the targets
        it would miss.
  \item We present \method, which ranks tiles with the VLM's own attention from a low-resolution view,
        re-scores only the uncertain tiles, and builds the KV cache by streaming, so that its cost is
        $\mathcal{O}(K)$ in the budget $K$ and independent of the number of tiles.
  \item On InternVL3.5 and five high-resolution benchmarks, \method keeps 95.0\% of the uncompressed
        accuracy with 5.0$\times$ less memory and is the most accurate method under every memory limit. On
        HR-Bench 8K, its peak memory stays at 0.41\,GB from 7 to 188 tiles, so one GPU serves at least 3.3$\times$ more
        requests than pruning inside the language model and 41$\times$ more than the uncompressed model
        (\cref{fig:scaling}).
\end{itemize}
